\chapter{Zaključak}
Elektronska trgovina predstavlja jedan od najznačajnijih domena primene informacionih sistema, pri čemu savremena rešenja dominantno počivaju na centralizovanom modelu poverenja. U ovoj disertaciji analizirana je mogućnost primene javnih blokčejn sistema kao alternativne infrastrukturne osnove za realizaciju elektronske trgovine, sa posebnim fokusom na uticaj decentralizacije na arhitekturu sistema, troškove izvršavanja i korisničko iskustvo.

Sprovedena teorijska analiza pokazala je da javni blokčejn sistemi sa podrškom za pametne ugovore mogu zadovoljiti referentni model elektronske trgovine, ali uz uvođenje dodatne složenosti koja proizlazi iz potrebe za decentralizovanom koordinacijom i upravljanjem zajedničkim stanjem sistema. Potvrđeno je da decentralizacija predstavlja dodatnu dimenziju projektovanja informacionih sistema, koja utiče na raspodelu kontrole, odgovornosti i poverenja među učesnicima, kao i na arhitekturu infrastrukturnih i tehnoloških resursa sistema.

Na konstruktivnom nivou dizajniran je arhitektonski model platforme za elektronsku trgovinu zasnovan na decentralizovanim \textit{Ethereum} i \textit{IPFS} tehnologijama. Predloženi model omogućava realizaciju poslovne logike korišćenjem pametnih ugovora, skladištenja podataka i korisnički kontrolisanog kriptografskog identiteta, čime se uklanja potreba za centralnim autoritetom kao nosiocem poverenja.

Eksperimentalna analiza pokazala je da troškovi izvršavanja osnovnih operacija nad entitetima elektronske trgovine dominantno zavise od broja pristupa trajnom skladištu podataka. Na osnovu toga razvijen je optimizovani model reprezentacije entiteta unutar jedne skladišne jedinice, zasnovan na proširenju postojećih NFT standarda. Rezultati komparativne analize standarda ERC-721, ERC-721A i ERC-721Psi potvrdili su da se primenom optimizovanih modela skladištenja mogu ostvariti značajna smanjenja potrošnje gasa prilikom izvršavanja operacija kreiranja, prenosa i uništavanja entiteta.

Dobijeni rezultati potvrđuju opštu hipotezu istraživanja da je moguće projektovati distribuiranu platformu elektronske trgovine zasnovanu na javnoj blokčejn mreži koja primenom optimizovanih modela skladištenja i obrade podataka omogućava smanjenje naknade transakcionih troškova korisnika uz očuvanje funkcionalnih i bezbednosnih svojstava sistema. Ključni izazov nije realizacija funkcionalnosti sistema, već efikasno upravljanje ograničenim računarskim i skladišnim resursima decentralizovane infrastrukture.

Istraživanje je pokazalo da optimizacija modela reprezentacije podataka i pristupa skladištu predstavlja jedan od ključnih faktora za smanjenje operativnih troškova decentralizovanih sistema, bez narušavanja njihovih osnovnih funkcionalnih svojstava. Istovremeno, primena decentralizovanog skladištenja podataka i prezentacionog sloja zasnovanih na IPFS mreži pokazuje da je moguće dodatno unaprediti stepen decentralizacije sistema uz očuvanje konzistentnosti i verifikabilnosti podataka.

Poseban značaj rezultata ogleda se u formalizaciji veze između decentralizacije sistema, načina organizacije podataka i ekonomskih troškova izvršavanja poslovnih procesa. Na taj način uspostavljen je metodološki okvir koji omogućava sistematsku analizu i projektovanje decentralizovanih informacionih sistema ne samo u domenu elektronske trgovine, već i u drugim oblastima primene blokčejn tehnologija.

Naučni doprinos ove disertacije ogleda se u formalizaciji referentnog modela elektronske trgovine u kontekstu javnih blokčejn sistema, definisanju analitičkog okvira za njihovu evaluaciju, razvoju optimizovanog modela reprezentacije entiteta unutar jedne skladišne jedinice, formalnom modelovanju mehanizama implicitnog vlasništva u NFT standardima, kao i eksperimentalnoj evaluaciji njihove troškovne efikasnosti.

Pored toga, predložen je arhitektonski model decentralizovane platforme za elektronsku trgovinu koji integriše javnu blokčejn mrežu, decentralizovano skladištenje podataka i decentralizovani prezentacioni sloj, čime je demonstrirana mogućnost realizacije potpuno decentralizovane arhitekture informacionog sistema elektronske trgovine.

Rezultati ovog istraživanja predstavljaju osnovu za dalji razvoj decentralizovanih informacionih sistema zasnovanih na javnim blokčejn mrežama. Buduća istraživanja mogu biti usmerena ka unapređenju interoperabilnosti između blokčejn sistema, optimizaciji performansi decentralizovanog skladištenja podataka, kao i razvoju novih modela za unapređenje privatnosti, skalabilnosti i efikasnosti decentralizovanih aplikacija.